\section{结论}

本文围绕一个核心问题展开：在尽量保留 Tucker 低秩骨架的前提下，如何以可控方式引入有效非线性。为此，我们提出了 NTD-PL。该方法以共享 Tucker 潜在张量作为结构主干，在条目层施加显式多项式链接，从而把非线性增强压缩到少量链接系数中。NTD-PL 更强调结构延续性、参数可控性与机制可解释性。

围绕这一建模思想，本文给出了从模型到实验的完整论证。理论上，NTD-PL 与标准 Tucker 构成清晰嵌套：线性特例可精确退化为 Tucker，次数增加形成受控扩张；在多项式生成情形下可实现精确匹配，在解析函数情形下可获得有限阶逼近并给出截断误差界。算法上，我们给出了统一的学习框架，可同时处理全观测分解与带掩码补全，并与实际实现保持一致。

实验结果支持了上述判断。在线性一致性实验中，受限 NTD-PL 与 Tucker 基本重合，说明模型不会在纯线性 regime 中制造伪增益；在可控非线性残差实验中，性能趋势与理论预期一致。真实高光谱实验进一步表明，在 CAVE 的全观测重构与随机缺失补全任务中，NTD-PL 在匹配参数预算下稳定优于 Tucker，且优势在 scene-wise 与 rate-wise 层面均具有一致性。

更重要的是，机制与空间证据解释了为什么有效。机制对照显示，Tucker 输出上的后处理多项式校准（PolyCal）只能部分逼近 NTD-PL，无法替代非线性链接与低秩骨架的联合建模。空间分析显示，收益并非全局均匀平滑，而主要集中在高难度且强边界的局部区域。阶数分析进一步表明，主要收益集中在低到中阶，边际增益随阶数上升逐步饱和，说明该方法的额外计算开销是可量化且有边界的。跨数据集结果也给出了清晰结论：NTD-PL 在多数真实场景中有效，但并非无条件成立。当数据更接近线性混合、条目层非线性偏离较弱时，其额外收益会明显减弱。

未来工作可沿三条路径推进：其一，研究更稳定或更灵活的显式链接族（如正交基或分段基）以改善高阶数值条件；其二，发展更高效的掩码感知训练与近似计算，提升大规模补全效率；其三，围绕可辨识性、泛化误差与阶数自适应选择开展更系统的理论分析。总体而言，NTD-PL 提供了一条可分析、可实现、可扩展的“增强 Tucker”路径，为结构化非线性低秩建模提供了可复用起点。
